\begin{register}{H}{LP\_GPIO\_CLK\_EN\_REG}{0x0000}\label{regdesc:LPGPIOCLKENREG}
\regfield{(reserved)}{31}{1}{0000000000000000000000000000000}%
\regfield{LP\_GPIO\_CLK\_EN}{1}{0}{{1}}%
\reglabel{Reset}\regnewline%
\begin{regdesc}\begin{reglist}
\label{fielddesc:LPGPIOCLKEN}\item [LP\_GPIO\_CLK\_EN] Configures whether or not to enable clock gate.\\
0: Not enable\\
1: Enable, the clock is free running. \\
(R/W)
\end{reglist}\end{regdesc}
\end{register}


\begin{register}{H}{LP\_GPIO\_VER\_DATE\_REG}{0x0004}\label{regdesc:LPGPIOVERDATEREG}
\regfield{(reserved)}{4}{28}{0000}%
\regfield{LP\_GPIO\_VER\_DATE}{28}{0}{{0x230323}}%
\reglabel{Reset}\regnewline%
\begin{regdesc}\begin{reglist}
\label{fielddesc:LPGPIOVERDATE}\item [LP\_GPIO\_VER\_DATE] Version control register. \\
(R/W)
\end{reglist}\end{regdesc}
\end{register}


\begin{register}{H}{LP\_GPIO\_OUT\_REG}{0x0008}\label{regdesc:LPGPIOOUTREG}
\regfield{(reserved)}{16}{16}{0000000000000000}%
\regfield{LP\_GPIO\_OUT\_DATA}{16}{0}{{0}}%
\reglabel{Reset}\regnewline%
\begin{regdesc}\begin{reglist}
\label{fielddesc:LPGPIOOUTDATA}\item [LP\_GPIO\_OUT\_DATA]  Configures the output value of GPIO0 \verb+~+ GPIO15 output in simple GPIO output mode.\\
0: Low level\\
1: High level\\
The value of bit0 \verb+~+ bit15 correspond to the output value of GPIO0 \verb+~+ GPIO15 respectively. Bit16 \verb+~+ bit31 are invalid.\\
(R/W/WTC)
\end{reglist}\end{regdesc}
\end{register}


\begin{register}{H}{LP\_GPIO\_OUT\_W1TS\_REG}{0x000C}\label{regdesc:LPGPIOOUTW1TSREG}
\regfield{(reserved)}{16}{16}{0000000000000000}%
\regfield{LP\_GPIO\_OUT\_DATA\_W1TS}{16}{0}{{0}}%
\reglabel{Reset}\regnewline%
\begin{regdesc}\begin{reglist}
\label{fielddesc:LPGPIOOUTDATAW1TS}\item [LP\_GPIO\_OUT\_DATA\_W1TS] Configures whether or not to set the output register \hyperref[regdesc:LPGPIOOUTREG]{LP\_GPIO\_OUT\_REG} of GPIO0 \verb+~+ GPIO15.\\
0: Not set\\
1: The corresponding bit in \hyperref[regdesc:LPGPIOOUTREG]{LP\_GPIO\_OUT\_REG} will be set to 1\\
Bit0 \verb+~+ bit15 are corresponding to GPIO0 \verb+~+ GPIO15. Bit16 \verb+~+ bit31 are invalid.
Recommended operation: use this register to set \hyperref[regdesc:LPGPIOOUTREG]{LP\_GPIO\_OUT\_REG}. \\
(WT)
\end{reglist}\end{regdesc}
\end{register}


\begin{register}{H}{LP\_GPIO\_OUT\_W1TC\_REG}{0x0010}\label{regdesc:LPGPIOOUTW1TCREG}
\regfield{(reserved)}{16}{16}{0000000000000000}%
\regfield{LP\_GPIO\_OUT\_DATA\_W1TC}{16}{0}{{0}}%
\reglabel{Reset}\regnewline%
\begin{regdesc}\begin{reglist}
\label{fielddesc:LPGPIOOUTDATAW1TC}\item [LP\_GPIO\_OUT\_DATA\_W1TC] Configures whether or not to clear the output register \hyperref[regdesc:LPGPIOOUTREG]{LP\_GPIO\_OUT\_REG} of GPIO0 \verb+~+ GPIO15 output.\\
0: Not clear\\
1: The corresponding bit in \hyperref[regdesc:LPGPIOOUTREG]{LP\_GPIO\_OUT\_REG} will be cleared.\\
Bit0 \verb+~+ bit15 are corresponding to GPIO0 \verb+~+ GPIO15. Bit16 \verb+~+ bit31 are invalid.
Recommended operation: use this register to clear \hyperref[regdesc:LPGPIOOUTREG]{LP\_GPIO\_OUT\_REG}. \\
(WT)
\end{reglist}\end{regdesc}
\end{register}


\begin{register}{H}{LP\_GPIO\_ENABLE\_REG}{0x0014}\label{regdesc:LPGPIOENABLEREG}
\regfield{(reserved)}{16}{16}{0000000000000000}%
\regfield{LP\_GPIO\_ENABLE\_DATA}{16}{0}{{0}}%
\reglabel{Reset}\regnewline%
\begin{regdesc}\begin{reglist}
\label{fielddesc:LPGPIOENABLEDATA}\item [LP\_GPIO\_ENABLE\_DATA]  Configures whether or not to enable the output of GPIO0 \verb+~+ GPIO15.\\
0: Not enable\\
1: Enable\\
Bit0 \verb+~+ bit15 are corresponding to GPIO0 \verb+~+ GPIO15. Bit16 \verb+~+ bit31 are invalid.\\
(R/W/WTC)
\end{reglist}\end{regdesc}
\end{register}


\begin{register}{H}{LP\_GPIO\_ENABLE\_W1TS\_REG}{0x0018}\label{regdesc:LPGPIOENABLEW1TSREG}
\regfield{(reserved)}{16}{16}{0000000000000000}%
\regfield{LP\_GPIO\_ENABLE\_DATA\_W1TS}{16}{0}{{0}}%
\reglabel{Reset}\regnewline%
\begin{regdesc}\begin{reglist}
\label{fielddesc:LPGPIOENABLEDATAW1TS}\item [LP\_GPIO\_ENABLE\_DATA\_W1TS] Configures whether or not to set the output enable register \hyperref[regdesc:LPGPIOENABLEREG]{LP\_GPIO\_ENABLE\_REG} of GPIO0 \verb+~+ GPIO15.\\
0: Not set\\
1: The corresponding bit in \hyperref[regdesc:LPGPIOENABLEREG]{LP\_GPIO\_ENABLE\_REG} will be set to 1\\
Bit0 \verb+~+ bit15 are corresponding to GPIO0 \verb+~+ GPIO15. Bit16 \verb+~+ bit31 are invalid. Recommended operation: use this register to set \hyperref[regdesc:LPGPIOENABLEREG]{LP\_GPIO\_ENABLE\_REG}.\\
(WT)
\end{reglist}\end{regdesc}
\end{register}


\begin{register}{H}{LP\_GPIO\_ENABLE\_W1TC\_REG}{0x001C}\label{regdesc:LPGPIOENABLEW1TCREG}
\regfield{(reserved)}{16}{16}{0000000000000000}%
\regfield{LP\_GPIO\_ENABLE\_DATA\_W1TC}{16}{0}{{0}}%
\reglabel{Reset}\regnewline%
\begin{regdesc}\begin{reglist}
\label{fielddesc:LPGPIOENABLEDATAW1TC}\item [LP\_GPIO\_ENABLE\_DATA\_W1TC] Configures whether or not to clear the output enable register \hyperref[regdesc:LPGPIOENABLEREG]{LP\_GPIO\_ENABLE\_REG} of GPIO0 \verb+~+ GPIO15. \\
0: Not clear\\
1: The corresponding bit in \hyperref[regdesc:LPGPIOENABLEREG]{LP\_GPIO\_ENABLE\_REG} will be cleared\\
Bit0 \verb+~+ bit15 are corresponding to GPIO0 \verb+~+ 15. Bit16 \verb+~+ bit31 are invalid. Recommended operation: use this register to clear \hyperref[regdesc:LPGPIOENABLEREG]{LP\_GPIO\_ENABLE\_REG}.\\
(WT)
\end{reglist}\end{regdesc}
\end{register}


\begin{register}{H}{LP\_GPIO\_STATUS\_REG}{0x0020}\label{regdesc:LPGPIOSTATUSREG}
\regfield{(reserved)}{16}{16}{0000000000000000}%
\regfield{LP\_GPIO\_STATUS\_DATA}{16}{0}{{0}}%
\reglabel{Reset}\regnewline%
\begin{regdesc}\begin{reglist}
\label{fielddesc:LPGPIOSTATUSDATA}\item [LP\_GPIO\_STATUS\_DATA]  The interrupt status of GPIO0 \verb+~+ GPIO15, can be configured by the software. \\
Bit0 \verb+~+ bit15 are corresponding to GPIO0 \verb+~+ GPIO15. Bit16 \verb+~+ bit31 are invalid. \\
Each bit represents the status of its corresponding GPIO:\\
0: Represents the GPIO does not generate the interrupt configured by \hyperref[fielddesc:LPGPIOPINNINTTYPE]{LP\_GPIO\_PIN\regindex{n}\_INT\_TYPE}, or this bit is configured to 0 by the software.\\
1: Represents the GPIO generates the interrupt configured by \hyperref[fielddesc:LPGPIOPINNINTTYPE]{LP\_GPIO\_PIN\regindex{n}\_INT\_TYPE}, or this bit is configured to 1 by the software.\\
(R/W/WTC)
\end{reglist}\end{regdesc}
\end{register}


\begin{register}{H}{LP\_GPIO\_STATUS\_W1TS\_REG}{0x0024}\label{regdesc:LPGPIOSTATUSW1TSREG}
\regfield{(reserved)}{16}{16}{0000000000000000}%
\regfield{LP\_GPIO\_STATUS\_DATA\_W1TS}{16}{0}{{0}}%
\reglabel{Reset}\regnewline%
\begin{regdesc}\begin{reglist}
\label{fielddesc:LPGPIOSTATUSDATAW1TS}\item [LP\_GPIO\_STATUS\_DATA\_W1TS] Configures whether or not to set the interrupt status register \hyperref[fielddesc:LPGPIOSTATUSDATA]{LP\_GPIO\_STATUS\_DATA} of GPIO0 \verb+~+ GPIO15. \\
Bit0 \verb+~+ bit15 are corresponding to GPIO0 \verb+~+ GPIO15. Bit16 \verb+~+ bit31 are invalid.\\
If the value 1 is written to a bit here, the corresponding bit in \hyperref[fielddesc:LPGPIOSTATUSDATA]{LP\_GPIO\_STATUS\_DATA} will be set to 1. \\
Recommended operation: use this register to set \hyperref[fielddesc:LPGPIOSTATUSDATA]{LP\_GPIO\_STATUS\_DATA}. \\
(WT)
\end{reglist}\end{regdesc}
\end{register}


\begin{register}{H}{LP\_GPIO\_STATUS\_W1TC\_REG}{0x0028}\label{regdesc:LPGPIOSTATUSW1TCREG}
\regfield{(reserved)}{16}{16}{0000000000000000}%
\regfield{LP\_GPIO\_STATUS\_DATA\_W1TC}{16}{0}{{0}}%
\reglabel{Reset}\regnewline%
\begin{regdesc}\begin{reglist}
\label{fielddesc:LPGPIOSTATUSDATAW1TC}\item [LP\_GPIO\_STATUS\_DATA\_W1TC]  Configures whether or not to clear the interrupt status register \hyperref[fielddesc:LPGPIOSTATUSDATA]{LP\_GPIO\_STATUS\_DATA} of GPIO0 \verb+~+ GPIO15. \\
Bit0 \verb+~+ bit15 are corresponding to GPIO0 \verb+~+ GPIO15. Bit16 \verb+~+ bit31 are invalid.\\
If the value 1 is written to a bit here, the corresponding bit in \hyperref[fielddesc:LPGPIOSTATUSDATA]{LP\_GPIO\_STATUS\_DATA} will be cleared. \\
Recommended operation: use this register to clear \hyperref[fielddesc:LPGPIOSTATUSDATA]{LP\_GPIO\_STATUS\_DATA}. \\(WT)
\end{reglist}\end{regdesc}
\end{register}


\begin{register}{H}{LP\_GPIO\_STATUS\_NEXT\_REG}{0x002C}\label{regdesc:LPGPIOSTATUSNEXTREG}
\regfield{(reserved)}{16}{16}{0000000000000000}%
\regfield{LP\_GPIO\_STATUS\_INTERRUPT\_NEXT}{16}{0}{{0x00}}%
\reglabel{Reset}\regnewline%
\begin{regdesc}\begin{reglist}
\label{fielddesc:LPGPIOSTATUSINTERRUPTNEXT}\item [LP\_GPIO\_STATUS\_INTERRUPT\_NEXT] Represents the interrupt source signal of GPIO0 \verb+~+ GPIO15.\\
Bit0 \verb+~+ bit15 are corresponding to GPIO0 \verb+~+ 15. Bit16 \verb+~+ bit31 are invalid. Each bit represents:\\
0: The GPIO does not generate the interrupt configured by \hyperref[fielddesc:LPGPIOPINNINTTYPE]{LP\_GPIO\_PIN\regindex{n}\_INT\_TYPE}.\\
1: The GPIO generates an interrupt configured by \hyperref[fielddesc:LPGPIOPINNINTTYPE]{LP\_GPIO\_PIN\regindex{n}\_INT\_TYPE}.\\
The interrupt could be rising edge interrupt, falling edge interrupt, level sensitive interrupt and any edge interrupt.\\  (RO)
\end{reglist}\end{regdesc}
\end{register}


\begin{register}{H}{LP\_GPIO\_IN\_REG}{0x0030}\label{regdesc:LPGPIOINREG}
\regfield{(reserved)}{16}{16}{0000000000000000}%
\regfield{LP\_GPIO\_IN\_DATA\_NEXT}{16}{0}{{0}}%
\reglabel{Reset}\regnewline%
\begin{regdesc}\begin{reglist}
\label{fielddesc:LPGPIOINDATANEXT}\item [LP\_GPIO\_IN\_DATA\_NEXT]  Represents the input value of GPIO0 \verb+~+ GPIO15. Each bit represents a pin input value:\\
0: Low level\\
1: High level\\
Bit0 \verb+~+ bit15 are corresponding to GPIO0 \verb+~+ GPIO15. Bit16 \verb+~+ bit31 are invalid.\\
(RO)
\end{reglist}\end{regdesc}
\end{register}


\begin{register}{H}{LP\_GPIO\_PIN\regindex{n}\_REG (\regindex{n}: 0 - 15)}{0x0034+4*\regindex{n}}\label{regdesc:LPGPIOPINNREG}
\regfield{(reserved)}{26}{6}{00000000000000000000000000}%
\regfield{LP\_GPIO\_PIN\regindex{n}\_EDGE\_WAKEUP\_CLR}{1}{5}{{0}}%
\regfield{LP\_GPIO\_PIN\regindex{n}\_PAD\_DRIVER}{1}{4}{{0}}%
\regfield{LP\_GPIO\_PIN\regindex{n}\_INT\_TYPE}{3}{1}{{0}}%
\regfield{LP\_GPIO\_PIN\regindex{n}\_WAKEUP\_ENABLE}{1}{0}{{0}}%
\reglabel{Reset}\regnewline%
\begin{regdesc}\begin{reglist}
\label{fielddesc:LPGPIOPINNWAKEUPENABLE}\item [LP\_GPIO\_PIN\regindex{n}\_WAKEUP\_ENABLE] Configures whether or not to enable GPIO wake-up function.\\
0: Disable\\
1: Enable\\
This function only wakes up system from Deep-sleep. \\(R/W)
\label{fielddesc:LPGPIOPINNINTTYPE}\item [LP\_GPIO\_PIN\regindex{n}\_INT\_TYPE]  Configures GPIO interrupt type.\\
0: GPIO interrupt disabled\\
1: Rising edge trigger\\
2: Falling edge trigger\\
3: Any edge trigger\\
4: Low level trigger\\
5: High level trigger\\ (R/W)
\label{fielddesc:LPGPIOPINNPADDRIVER}\item [LP\_GPIO\_PIN\regindex{n}\_PAD\_DRIVER] Configures to select pin drive mode. \\
0: Normal output\\
1: Open drain output \\(R/W)
\label{fielddesc:LPGPIOPINNEDGEWAKEUPCLR}\item [LP\_GPIO\_PIN\regindex{n}\_EDGE\_WAKEUP\_CLR]  Configures to clear edge-triggered GPIO wake-up source in PMU. \\
1: The corresponding bit in LP\_PMU\_xxx will be cleared.\\
0: No effect.\\
(WT)
\end{reglist}\end{regdesc}
\end{register}


\begin{register}{H}{LP\_GPIO\_FUNC\regindex{n}\_IN\_SEL\_CFG\_REG (\regindex{n}: 0 - 15)}{0x0074+4*\regindex{n}}\label{regdesc:LPGPIOFUNCNINSELCFGREG}
\regfield{(reserved)}{24}{8}{000000000000000000000000}%
\regfield{LP\_GPIO\_FUNC\regindex{n}\_IN\_SEL}{6}{2}{{0x30}}%
\regfield{LP\_GPIO\_SIG\regindex{n}\_IN\_SEL}{1}{1}{{0}}%
\regfield{LP\_GPIO\_FUNC\regindex{n}\_IN\_INV\_SEL}{1}{0}{{0}}%
\reglabel{Reset}\regnewline%
\begin{regdesc}\begin{reglist}
\label{fielddesc:LPGPIOFUNCNININVSEL}\item [LP\_GPIO\_FUNC\regindex{n}\_IN\_INV\_SEL]  Configures whether or not to invert the input value.\\
0: Not invert\\
1: Invert\\ (R/W)
\label{fielddesc:LPGPIOSIGNINSEL}\item [LP\_GPIO\_SIG\regindex{n}\_IN\_SEL]  Configures whether or not to route signals via LP GPIO matrix.\\
0: Bypass LP GPIO matrix, i.e., connect signals directly to peripheral configured in LP IO MUX.\\
1: Route signals via LP GPIO matrix.\\  (R/W)
\label{fielddesc:LPGPIOFUNCNINSEL}\item [LP\_GPIO\_FUNC\regindex{n}\_IN\_SEL]  Configures to select a pin from the 16 LP GPIO pins to connect the input signal \regindex{n}.\\
0: Select GPIO0\\
1: Select GPIO1\\
......\\
14: Select GPIO14\\
15: Select GPIO15\\
16 \verb+~+ 31: Invalid \\
32 \verb+~+ 47: A constantly low input\\
48 \verb+~+ 63: A constantly high input\\
(R/W)
\end{reglist}\end{regdesc}
\end{register}

\begin{register}{H}{LP\_GPIO\_FUNC\regindex{n}\_OUT\_SEL\_CFG\_REG (\regindex{n}: 0 - 15)}{0x00F4+4*\regindex{n}}\label{regdesc:LPGPIOFUNCNOUTSELCFGREG}
\regfield{(reserved)}{23}{9}{00000000000000000000000}%
\regfield{LP\_GPIO\_FUNC\regindex{n}\_OUT\_SEL}{6}{3}{{0x20}}%
\regfield{LP\_GPIO\_FUNC\regindex{n}\_OUT\_INV\_SEL}{1}{2}{{0}}%
\regfield{LP\_GPIO\_FUNC\regindex{n}\_OE\_SEL}{1}{1}{{0}}%
\regfield{LP\_GPIO\_FUNC\regindex{n}\_OE\_INV\_SEL}{1}{0}{{0}}%
\reglabel{Reset}\regnewline%
\begin{regdesc}\begin{reglist}
\label{fielddesc:LPGPIOFUNCNOEINVSEL}\item [LP\_GPIO\_FUNC\regindex{n}\_OE\_INV\_SEL]  Configures whether or not to invert the output enable signal.\\
0: Not invert\\
1: Invert\\ (R/W)
\label{fielddesc:LPGPIOFUNCNOESEL}\item [LP\_GPIO\_FUNC\regindex{n}\_OE\_SEL]  Configures to select the source of output enable signal.\\
0: Use output enable signal from peripheral.\\
1: Force the output enable signal to be sourced from bit\regindex{n} of \hyperref[regdesc:LPGPIOENABLEREG]{LP\_GPIO\_ENABLE\_REG}. \\(R/W)
\label{fielddesc:LPGPIOFUNC0OUTINVSEL}\item [LP\_GPIO\_FUNC\regindex{n}\_OUT\_INV\_SEL] Configures whether or not to invert the output value.\\
0: Not invert\\
1: Invert\\ (R/W)
\label{fielddesc:LPGPIOFUNCNOUTSEL}\item [LP\_GPIO\_FUNC\regindex{n}\_OUT\_SEL] Configures to select a signal \regindex{Y} (0 <= \regindex{Y} < 32) from 31 peripheral signals to be output to GPIO\regindex{n}.\\
0: Select signal 0\\
1: Select signal 1\\
......\\
30: Select signal 30\\
31: Select signal 31\\
Or\\
32: Bit\regindex{n} of \hyperref[regdesc:LPGPIOOUTREG]{LP\_GPIO\_OUT\_REG} and \hyperref[regdesc:LPGPIOENABLEREG]{LP\_GPIO\_ENABLE\_REG} are selected as the output value and output enable.

For the detailed signal list, see Table \ref{tab:iomuxgpio-periph-signals-via-lp-gpio-matrix-en}. \\(R/W)
\end{reglist}\end{regdesc}
\end{register}

